% Generated draft: requires mathematical and source review.
\documentclass[11pt]{article}
\usepackage[margin=25mm]{geometry}
\usepackage{fontspec}
\setmainfont{Latin Modern Roman}
\usepackage{amsmath,amssymb,mathtools,unicode-math}
\setmathfont{Latin Modern Math}
\usepackage{enumitem,needspace,xurl,hyperref}
\setlength{\parindent}{0pt}
\setlength{\parskip}{5pt}
\setlength{\emergencystretch}{3em}
\newcommand{\N}{\mathbb{N}}
\newcommand{\Z}{\mathbb{Z}}
\newcommand{\Q}{\mathbb{Q}}
\newcommand{\R}{\mathbb{R}}
\newcommand{\C}{\mathbb{C}}
\newcommand{\F}{\mathbb{F}}
\newcommand{\A}{\mathbb{A}}
\newcommand{\PP}{\mathbb{P}}
\newcommand{\E}{\mathbb{E}}
\newcommand{\eps}{\varepsilon}
\newcommand{\dd}{\,\mathrm{d}}
\begin{document}
\section*{UnsolvedMath ID: 800260}
\textbf{Source problem code:} catalogue entry
\subsection{Definitions and mathematical statement}
Let \(M\) be a Stein manifold of complex dimension \(n\), and let \(T^{1,0}M\) denote its holomorphic tangent bundle. Say that \(M\) is holomorphically parallelizable if
\[
T^{1,0}M\cong M\times\C^n
\]
as holomorphic vector bundles. A holomorphic map \(F:M\to\C^n\) is an immersion if its differential
\[
\dd F_p:T^{1,0}_pM\longrightarrow\C^n
\]
is injective for every \(p\in M\). Since both spaces have dimension \(n\), this means that \(\dd F_p\) is an isomorphism at every point; global injectivity of \(F\) is not required.

Under the holomorphic interpretation of parallelizability, the question is: does every holomorphically parallelizable Stein manifold \(M\) of complex dimension \(n\) admit a holomorphic immersion into \(\C^n\)? Equivalently, must there exist holomorphic functions \(f_1,\ldots,f_n\) on \(M\) such that
\[
\dd f_1\wedge\cdots\wedge\dd f_n\ne0
\quad\text{at every point of }M?
\]

\textbf{Unresolved qualification.} The source says only ``parallelizable'' and does not specify the category. If smooth real parallelizability is intended, the hypothesis instead means that the underlying real tangent bundle satisfies \(T_{\R}M\cong M\times\R^{2n}\) as smooth real vector bundles, with the same requested holomorphic immersion. These hypotheses must not be silently identified. Dimension is interpreted here as complex dimension, as suggested by the Stein setting and target \(\C^n\).

The supplied source states that the answer is affirmative for \(n=1\), attributing this to Gunning and Narasimhan, and that such an \(M\) admits an immersion into \(\C^{n+1}\), attributing this to Eliashberg--Gromov. These are retained as contextual claims of the source, with the same unresolved meaning of ``parallelizable''; they have not been independently verified here.
\subsection{Short English statement}
The question asks whether every parallelizable Stein manifold of complex dimension \(n\) admits a holomorphic map into \(\C^n\) whose differential is invertible everywhere. The source does not specify whether parallelizability is holomorphic or smooth real. Its positive statements for dimension one and for target \(\C^{n+1}\) are supplied background claims.
\subsection{Sources}
\begin{itemize}
\item Local UnsolvedMath registry, permanent ID \(800260\), exact source code \texttt{catalogue entry}, title ``Several Complex Variables Problem List --- Problem 59.'' The supplied record labels the problem \texttt{open}; this is a registry label, not an independently verified current status. The record was updated on 2026-10-16.
\item Buzzard, Fornaess, and Krantz, \emph{The Several Complex Variables Problem List} (1995), Problem 59. Supplied source URL: \url{https://arxiv.org/abs/math/9509202}. The supplied provenance records access on 2026-10-16 and extraction from arXiv LaTeX; no source visit is claimed here.
\item The problem text is signed R. Narasimhan, Department of Mathematics, University of Chicago, Chicago, IL 60637. The registry's \texttt{proposed\_by} field names Buzzard, Fornaess, and Krantz, and its proposed year is 1995; the list attribution and the signed contributor attribution are both preserved.
\item The source attributes the dimension-one result to Gunning and Narasimhan and the \(\C^{n+1}\) immersion result to Eliashberg--Gromov. No individual bibliographic references or links for those results are supplied.
\end{itemize}
\end{document}
